\section{HIPÓTESIS}

\subsection{Planteamiento de la Hipótesis}

Para evaluar emp\'{i}ricamente la viabilidad del protocolo dise\~{n}ado bajo condiciones de estr\'{e}s transaccional, se definen las siguientes hip\'{o}tesis de investigaci\'{o}n:

\begin{itemize}
    \item \textbf{Hip\'{o}tesis Alternativa ($H_1$):} La integraci\'{o}n de un protocolo de votaci\'{o}n electr\'{o}nica con criptograf\'{i}a post-cu\'{a}ntica basada en ret\'{i}culos modulares (ML-KEM y ML-DSA) y privacidad algor\'{i}tmica (zk-SNARKs Groth16) sobre un entorno de contratos inteligentes, permite obtener una latencia promedio de procesamiento por voto menor o igual a 200 milisegundos sobre hardware convencional, logrando as\'{i} la viabilidad operativa del sistema.
    \item \textbf{Hip\'{o}tesis Nula ($H_0$):} La integraci\'{o}n de un protocolo de votaci\'{o}n electr\'{o}nica con criptograf\'{i}a post-cu\'{a}ntica basada en ret\'{i}culos modulares (ML-KEM y ML-DSA) y privacidad algor\'{i}tmica (zk-SNARKs Groth16) sobre un entorno de contratos inteligentes, genera una latencia promedio de procesamiento por voto superior a 200 milisegundos sobre hardware convencional, imposibilitando la viabilidad operativa del sistema.
\end{itemize}

\subsection{Definición y Clasificación de Variables}

La evaluaci\'{o}n experimental del protocolo propuesto requiere el control y medici\'{o}n de las siguientes variables clasificadas de manera jer\'{a}rquica:

\begin{itemize}
    \item \textbf{Variable Independiente (Causa):}
    \begin{itemize}
        \item \textit{Protocolo criptogr\'{a}fico h\'{i}brido poscu\'{a}ntico y de privacidad:} Definido por la implementaci\'{o}n conjunta de algoritmos de cifrado y firma resistentes a la computaci\'{o}n cu\'{a}ntica (est\'{a}ndares FIPS~203 y FIPS~204) e integraci\'{o}n de circuitos aritm\'{e}ticos de conocimiento cero (zk-SNARK Groth16).
    \end{itemize}
    
    \item \textbf{Variable Dependiente (Efecto):}
    \begin{itemize}
        \item \textit{Viabilidad y eficiencia operativa del sistema de votaci\'{o}n:} Definida como la capacidad de procesamiento del protocolo en t\'{e}rminos de velocidad transaccional, consumo de recursos f\'{i}sicos de infraestructura y costos de ejecuci\'{o}n de l\'{o}gica on-chain.
    \end{itemize}

    \item \textbf{Variables de Control (Constantes):}
    \begin{itemize}
        \item \textit{Rigor matem\'{a}tico del est\'{a}ndar criptogr\'{a}fico:} El cumplimiento obligatorio de los vectores de prueba oficiales emitidos por el NIST (CAVP) para validar la correcci\'{o}n de la ejecuci\'{o}n criptogr\'{a}fica.
        \item \textit{Infraestructura f\'{i}sica del entorno de simulaci\'{o}n:} Las especificaciones f\'{i}sicas constantes del servidor local donde se realiza la simulaci\'{o}n de carga masiva (CPU, memoria RAM y velocidad de almacenamiento).
    \end{itemize}
\end{itemize}

\subsection{Justificaci\'{o}n de Exclusiones en la Matriz de Variables}

Con el fin de preservar el rigor del dise\~{n}o experimental y evitar el principio de multidimensionalidad confusa, se aclara la naturaleza metodol\'{o}gica de dos componentes fundamentales del marco te\'{o}rico que no figuran como variables operacionales en la matriz:

\begin{itemize}
    \item \textbf{Computaci\'{o}n Cu\'{a}ntica:} No constituye una variable independiente ni dependiente debido a que la investigaci\'{o}n no manipula ni mide capacidades de procesamiento cu\'{a}ntico real. Su rol metodol\'{o}gico se limita a ser una \textit{Variable de Entorno y Motivaci\'{o}n Est\'{a}tica}. El vector de amenaza cu\'{a}ntica es mitigado de manera constante mediante la adopci\'{o}n de los est\'{a}ndares matem\'{a}ticos normados por el NIST (FIPS~203 y FIPS~204), cuyos par\'{a}metros ingresan al experimento de forma determinista y fija.
    \item \textbf{Pruebas de Conocimiento Cero ($zk$-SNARKs):} No se operacionalizan como variables debido a que el protocolo de privacidad (Groth16 implementado en Circom\footnote{Circom es un lenguaje y compilador dise\~{n}ado para escribir circuitos aritm\'{e}ticos que representan declaraciones l\'{o}gicas, mientras que snarkjs es la biblioteca de soporte para generar y verificar las pruebas matem\'{a}ticas de conocimiento cero correspondientes.}) act\'{u}a como una \textit{Constante de Dise\~{n}o Arquitect\'{o}nico}. El circuito aritm\'{e}tico y sus restricciones matem\'{a}ticas permanecen invariables e id\'{e}nticos durante la ejecuci\'{o}n de los $N = 10{,}000$ eventos de sufragio masivo. Su presencia es un requisito mandatorio de seguridad y no un elemento sujeto a gradaci\'{o}n o comparaci\'{o}n experimental.
\end{itemize}

Al aislar estos componentes como constantes y factores de control, el modelo experimental garantiza que cualquier fluctuaci\'{o}n registrada en la latencia temporal, consumo de CPU/RAM o costos de gas sea atribuible \'{u}nica y exclusivamente a la carga operativa de las firmas basadas en ret\'{i}culos bajo escenarios de alta concurrencia.

\subsection{Operacionalización de Variables}

Con el fin de cuantificar de forma objetiva las variables planteadas, se establece la matriz de operacionalizaci\'{o}n en la Tabla \ref{tab:operacionalizacion}.

\begin{table}[H]
\centering
\caption{Matriz de operacionalizaci\'{o}n de variables.}
\label{tab:operacionalizacion}
\renewcommand{\arraystretch}{1.3}
\fontsize{6.5}{7.8}\selectfont
\begin{tabular}{|p{1.4cm}|p{1.0cm}|p{1.9cm}|p{1.9cm}|p{1.5cm}|p{1.8cm}|p{1.3cm}|p{1.6cm}|}
\hline
\textbf{Variable} & \textbf{Tipo} & \textbf{Definición Conceptual} & \textbf{Definición Operacional} & \textbf{Dimensión} & \textbf{Indicadores} & \textbf{Unidad / Valor} & \textbf{Instrumento / Fuente} \\
\hline
\multirow{2}{1.4cm}{Protocolo criptogr\'{a}fico h\'{i}brido poscu\'{a}ntico y de privacidad} & \multirow{2}{1.0cm}{Independiente} & Algoritmos criptogr\'{a}ficos basados en ret\'{i}culos modulares para seguridad cu\'{a}ntica combinados con pruebas de anonimato. & Implementaci\'{o}n del esquema de cifrado ML-KEM-768 y firma digital ML-DSA-65 en software cliente. & Criptograf\'{i}a de ret\'{i}culos (PQC) & Implementaci\'{o}n de est\'{a}ndares FIPS~203 y FIPS~204 & Cumplimiento de especificaciones & C\'{o}digo del protocolo / Suite de pruebas \\
\cline{3-8}
 &  & Criptograf\'{i}a de conocimiento cero para validar la informaci\'{o}n del sufragio de manera an\'{o}nima. & Construcci\'{o}n del circuito de verificaci\'{o}n en lenguaje Circom aplicando Groth16 y hash Poseidon. & Privacidad algor\'{i}tmica (ZKP) & Complejidad del circuito aritm\'{e}tico & N\'{u}mero de restricciones (constraints) & Compilador de Circom / snarkjs \\
\hline
\multirow{3}{1.4cm}{Viabilidad y eficiencia operativa del sistema de votaci\'{o}n} & \multirow{3}{1.0cm}{Dependiente} & Capacidad de respuesta del protocolo en relaci\'{o}n al tiempo de atenci\'{o}n de los registros. & Medici\'{o}n del tiempo transcurrido desde el env\'{i}o de la transacci\'{o}n hasta la recepci\'{o}n de respuesta. & Rendimiento temporal & Latencia total por sufragio & Milisegundos (ms) [Meta: $\le$ 200 ms] & M\'{o}dulo de medici\'{o}n de tiempo / script Python \\
\cline{3-8}
 &  & Consumo de capacidad f\'{i}sica de c\'{o}mputo necesario para sostener la operaci\'{o}n. & Registro de la carga de procesador y uso de memoria vol\'{a}til en los nodos de red durante pruebas. & Consumo de infraestructura & Porcentaje de uso de CPU y volumen de RAM & CPU (\%) y RAM (MB/GB) & Monitor de recursos del S.O. / scripts de diagn\'{o}stico \\
\cline{3-8}
 &  & Gasto de procesamiento l\'{o}gico tasado en la cadena de bloques por ejecuci\'{o}n. & Cantidad de recursos consumidos on-chain para verificar las firmas y pruebas en Solidity. & Costo de red & Consumo transaccional on-chain & Unidades de gas consumidas en la EVM & Hardhat Gas Reporter / nodo de red \\
\hline
Rigor matem\'{a}tico del est\'{a}ndar criptogr\'{a}fico & Control & Veracidad del algoritmo respecto al comportamiento matem\'{a}tico oficial definido. & Verificaci\'{o}n de la salida del software utilizando los conjuntos de prueba del NIST. & Rigor del est\'{a}ndar & Cumplimiento de vectores CAVP & Valor booleano (Cumple / No cumple) & Vectores oficiales de prueba CAVP / pytest \\
\hline
Infraestructura del entorno de simulaci\'{o}n & Control & Conjunto de propiedades de hardware y red estables asignadas al experimento. & Delimitaci\'{o}n de recursos de hardware fijos durante las pruebas en el entorno de desarrollo. & Entorno f\'{i}sico & Especificaciones del servidor convencional & CPU, RAM y almacenamiento fijos & Ficha t\'{e}cnica de hardware / informaci\'{o}n del sistema \\
\hline
\end{tabular}
\par\vspace{3pt}{\footnotesize\textit{Fuente: Elaboraci\'{o}n propia.}}
\end{table}

